\section{Main Law}

\subsection{Problem statement}

Consider a one-dimensional observation stream
\[
  Y_t = \mu_t + \varepsilon_t,
  \qquad
  \varepsilon_t \stackrel{\mathrm{iid}}{\sim} \mathcal N(0,\sigma^2),
\]
where the mean path $\mu_t$ evolves under a local H\"older envelope with
exponent $H\in(0,1]$ and roughness scale $\zeta_t$. For a fixed operating
memory level $n$, the temporal-validity horizon from finite-memory tracking
provides a local benchmark scale
\[
  n_t^* \asymp (C_K / \zeta_t)^{1/(a+H)}
\]
with $a=1/2$ in the Gaussian root-$n$ benchmark regime.

Define the validity-expiry time of the operating memory level by
\[
  \tau_{\mathrm{valid}}(n)
  \coloneqq
  \inf\{t \ge 1 : n > n_t^*\},
\]
or its persistent analogue when one wants to suppress transient crossings.
Define $\tau_{\mathrm{detect}}$ as the stopping time of a sequential detector
run on the same filtration. The central delay object is
\[
  \Delta_{\mathrm{val-det}}(n)
  \coloneqq
  \tau_{\mathrm{detect}} - \tau_{\mathrm{valid}}(n).
\]

The detector class is indexed by a finite-horizon false-alarm budget.
Let $\mathbb P_0$ denote the stationary Gaussian no-change law with the same
noise variance and a constant mean path. For a horizon $T \in \mathbb N$ and a
budget $\alpha \in (0,1)$, define
\[
  \mathcal D_{\alpha,T}
  \coloneqq
  \bigl\{\tau : \tau \text{ is a stopping time and } \mathbb P_0(\tau \le T) \le \alpha\bigr\}.
\]
This is the exact false-alarm notion used here: under the stationary null, the
probability of at least one alarm before the calibration horizon $T$ is at most
$\alpha$.

The benchmark question is whether, under that constraint, the detector can fire
before validity expiry on a drifting Gaussian location path. The target is a
positive lower bound on the sign of $\Delta_{\mathrm{val-det}}(n)$.

\subsection{Detector class and false-alarm notion}

The detector class is deliberately broad. It contains any adapted stopping time
calibrated to a null false-alarm budget on a fixed horizon. Standard detectors
such as ADWIN, Page--Hinkley, CUSUM, and KSWIN enter only as benchmark
instantiations once they are tuned to satisfy the same $\mathbb P_0(\tau \le T)
\le \alpha$ constraint.

The finite-horizon choice matches the experimental calibration. It is also the
cleanest notion for the first theorem because it separates the alarm budget from
the drift geometry. A sharper average-run-length formulation can be studied
later, but it is not needed for the benchmark sign result.

\subsection{Candidate delay law}

\paragraph{Benchmark drift family.}
Let $t_\star$ denote the validity-expiry time of the operating window, and
consider the one-sided H\"older location path
\[
  \mu_t =
  \begin{cases}
    0, & t \le t_\star, \\
    \zeta (t - t_\star)^H, & t > t_\star.
  \end{cases}
\]
This path is stationary before expiry and drifts only after the temporal-validity
clock has crossed the chosen memory level.

\paragraph{Theorem (first benchmark sign result).}
Fix $\alpha \in (0,1)$ and $T \ge t_\star$. For every stopping time
$\tau \in \mathcal D_{\alpha,T}$ and every benchmark path of the form above,
\[
  \mathbb P_\mu\!\left(\tau > \tau_{\mathrm{valid}}(n)\right)
  \ge 1 - \alpha.
\]
Equivalently, the validity-detection delay satisfies
\[
  \mathbb P_\mu\!\left(\Delta_{\mathrm{val-det}}(n) > 0\right)
  \ge 1 - \alpha.
\]
The proof is immediate from the definition of $\mathcal D_{\alpha,T}$ because
the drifting benchmark path agrees with the null law on $\{1,\dots,t_\star\}$.
The sign result is therefore a benchmark consequence of the false-alarm budget
and the null prefix, not a detector-specific property.

\paragraph{Stronger target (delay lower bound).}
The sharper theorem target is a quantitative lower bound
\[
  \inf_{\tau \in \mathcal D_\alpha}
  \sup_{\mu \in \mathcal M(H,\zeta)}
  \mathbb E_\mu\!\left[(\tau - \tau_{\mathrm{valid}}(n))_+\right]
  \ge c_1 \, g(\alpha,\zeta,H,n),
\]
for some explicit positive function $g$ on the pre-detection invalidity regime.
The benchmark model should reveal how $g$ depends on the local evidence growth
after horizon crossing. The current theorem only isolates the sign; the lower
bound is the open next step.

\paragraph{Mechanism.}
The proof route separates two scales. The validity clock is determined by the
finite-memory horizon, which contracts once the local roughness level rises. The
detection clock is determined by how quickly post-expiry evidence can accumulate
without violating the null false-alarm constraint. In the benchmark sign result,
the null prefix makes the delay sign a direct consequence of the calibration
budget. In the stronger lower-bound target, the same separation should be proved
without an exact null prefix, using contiguity over the post-expiry window.

\paragraph{Benchmark theorem program.}
The tractable benchmark is a Gaussian location path whose local increments follow
a one-sided H\"older profile. The intended theorem line has four steps.

1. Prove a finite-memory expiry bound for $\tau_{\mathrm{valid}}(n)$ from the
   closed horizon law.
2. Bound the post-expiry information accumulated by time
   $\tau_{\mathrm{valid}}(n)+m$.
3. Convert the null false-alarm constraint into a lower bound on required alarm
   evidence for the detector class.
4. Deduce a sign law, and then a quantitative lower bound, for
   $\Delta_{\mathrm{val-det}}(n)$.

At present, the exact null-prefix benchmark theorem is the first theorem-level
result in the line. The stronger contiguity-based lower bound remains open. The
existing false-alarm-calibrated experiments support the same sign pattern across
multiple detector families and across $H\in\{0.5,0.75,1.0\}$, but theorem-level
closure for the quantitative delay law remains to be proved.
